\chapter{Conclusions}\label{ch:conclusions}
\section{Summary of the Work}
The project set out to turn a photograph-level ROP classifier into an eye-level screening system with a different model, higher accuracy and a publishable evaluation. We fixed the evaluation first: patient-grouped five-fold cross-validation over 6{,}700 photographs from three hospitals, a fourth hospital held out, one pre-registered primary endpoint, a selection rule that never looks at the held-out hospital, clustered bootstrap intervals and a stated minimum detectable effect. We then built the eye-level route and PDF report in the application, implemented a training pipeline for vision transformers on free Kaggle GPUs, and trained and scored seven models with two seeds each: a retina-pretrained transformer, a natural-image-pretrained transformer of identical architecture, the EfficientNet baseline, an examination-level training variant, an ensemble, an adaptation variant and a native-resolution variant. A final pre-registered round audited the training data (per-camera frame quality and a pixel-checked near-duplicate sweep) and added two more models, the selected transformer on the cleaned pool and the same pool with a regularising recipe, judged against a pre-set adoption margin and a leave-one-hospital-out guard. Every run, prediction, score and figure is committed and reproducible by one command.

\section{Conclusions against Objectives}
\begin{enumerate}
\item \textbf{O1, eye-level inference: achieved.} The application takes every photograph of one eye and returns one stage by the worst-view rule, the driving photograph and a PDF; the route and its error cases are covered by automated tests.
\item \textbf{O2, a different backbone, compared fairly: achieved.} DINOv2 ViT-S/14, selected by the pre-registered rule with a cross-validation QWK of 0.825 (0.737--0.911), replaces EfficientNet in the application, served as the mean of its five fold models (\DeployedQWK{} per eye at the new hospital). The retina-pretrained RETFound-Green scored 0.794 in cross-validation and 0.774 per eye at the new hospital; the difference between the two pre-training domains is $+0.018$ (0.000--0.041) per eye in favour of the retinal model and $-0.031$ ($-0.077$ to $+0.015$) in cross-validation: no consistent advantage either way.
\item \textbf{O3, pre-registered evaluation: achieved.} Protocol v2 and five dated amendments were committed before the experiments they govern; the v4 adoption rule kept model~B, as written, although model~I scored better at the new hospital; the locked half remains unopened for these models.
\item \textbf{O4, accuracy targets: achieved for the stated thresholds, not beyond the baseline.} Cross-validation QWK 0.825 against the 0.80 target and eye-level QWK 0.756 against 0.75; accuracy 76.6\% per photograph and 63.0\% per eye; any-ROP AUC 0.94 per eye. No model beats the baseline's eye-level QWK of 0.760 by more than its interval (the best, model~I, scores 0.776), so the review's request for higher accuracy on every metric was not met, and we report that the ceiling lies in the labels rather than the model.
\item \textbf{O5, pre-registered improvements: achieved, with mostly null results.} Resolution (518~px), ensembling, ordinal decoding and unsupervised adaptation did not improve the eye-level result, and neither the cleaned training pool nor the regularising recipe reached the $+0.010$ adoption margin in cross-validation (H $-0.013$, I $-0.004$). Model~I's gain at the new hospital ($+0.020$, $+0.003$ to $+0.041$ over B) is descriptive and below the detectable difference. Two gains are real: stage-1 eye recall at the new hospital rose from 0.40 to 0.50--0.52, and training at the examination unit raised Ostrava examination-level QWK from 0.701 to 0.846 ($+0.145$, 0.094--0.257, post hoc).
\item \textbf{O6, manuscript: partly achieved.} The draft, with every number traced, is written; it awaits the single approved scoring of the locked half for its primary endpoint and the author details for submission.
\end{enumerate}

\section{Scope for Future Work}
\begin{enumerate}
\item \textbf{Score the locked half once}, after the guide's approval and the log-book entry, for all nine models together; this gives the pre-registered primary endpoint and completes the manuscripts. It cannot confirm a gain as small as model~I's, but it will show whether the ordering holds on unseen eyes.
\item \textbf{Re-grade the stage-1 boundary with a clinician.} The nine models make the same adjacent-stage errors, chiefly stage~1 called normal; re-grading the stage-1 photographs on which several models disagree with the label, together with the 14 near-duplicate pairs that carry different labels, would measure the label noise directly and is the one step that can raise the ceiling.
\item \textbf{Set the screening threshold per site.} The cross-validation threshold transfers at 64--68\% sensitivity; a small labelled calibration set from each new hospital would restore the 90\% target.
\item \textbf{Extend examination-level training} to any cohort with session or eye labels, and test it on a second hospital with eye-level labels; the Ostrava gain is the one result worth building on. More infants with stage~4/5 and AP-ROP from more than one source are the data, not the recipe, that the leave-one-hospital-out folds show to be missing.
\item \textbf{Clinical validation}: a prospective study with a paediatric ophthalmologist as the reference standard, which is outside what a student team can do but is the precondition for any use on patients.
\item \textbf{Application}: camera integration, plus-disease and zone outputs when consistently labelled data exist, and a model card stating the licences and the site-validation requirement.
\end{enumerate}
